\title{Literature Review: Deep Learning for Seizure Detection and Prediction
from Non-EEG Wearable Signals (2015--2025)}
\author{Paulin Saher}
\date{\today}
\maketitle

\section{Introduction}
\subsection{Problem Statement and Motivation}
Artificial intelligence (AI), particularly deep learning for time-series
data, has enabled substantial advances in medical signal analysis,
supporting earlier detection and continuous monitoring. Epilepsy affects
roughly 7.6 per 1000 people and about 30\% of patients are drug
resistant, motivating demand for timely, low-burden monitoring and early
intervention \parencite{beghiEpidemiologyEpilepsy2019,kwanEarlyIdentificationRefractory2000,chenNewEraPersonalised2020}.
While EEG remains the reference standard in controlled clinical settings,
its obtrusiveness limits pervasive ambulatory use. Wearable sensing of
autonomic and movement signals (ECG/HRV, PPG, EDA, ACC) offers the
potential for continuous monitoring outside the clinic.

Healthcare constraints sharpen modeling objectives: (i) achieving high
sensitivity while controlling false alarms per hour (FAR/h) to avoid
alarm fatigue; (ii) balancing performance against interpretability and
safety requirements for regulated contexts; (iii) ensuring robustness to
domain shift across patients, devices, sites, and activities; and (iv)
addressing data scarcity in prospective, multimodal, patient-preserving
designs. The seminar theme ("Trends in DL / Time Series
Transformers") is operationalized here via a focused medical use case:
comparing traditional feature-based models and emerging deep temporal
architectures (1D-CNN, LSTM, TCN, attention/transformers) for seizure
detection and preictal prediction from non-EEG wearable biosignals.

The key question is which modeling and evaluation choices deliver
clinically viable, low FAR/h, high-sensitivity performance under
pseudo-/prospective, ambulatory conditions.

\subsection{Aim and Research Questions}
This literature review synthesizes machine and deep learning approaches
for seizure detection and preictal prediction using non-EEG
wearable/autonomic biosignals (ECG/HRV, PPG, EDA, ACC) reported between
2015 and 2025, to identify design and evaluation strategies associated
with clinically meaningful ambulatory performance.

\paragraph{Primary Research Question}
Among non-EEG wearable biosignals (ECG/HRV, PPG, EDA, ACC) from
2015--2025, which modeling and evaluation choices yield low alarm-based
FAR/h with high sensitivity and acceptable latency in ambulatory
seizure detection and preictal prediction under pseudo-/prospective
designs?

\paragraph{Sub-questions}
\begin{enumerate}
  \item What ranges of sensitivity, FAR/h, latency (and secondary
    metrics: AUC, F1) are reported across retrospective,
    pseudo-prospective, and prospective designs, and how do alarm-based
    results differ from sample-based metrics?
  \item How do traditional feature-based models compare to deep
    temporal architectures (1D-CNN, TCN, LSTM) and emerging
    attention/transformer variants regarding alarm-based performance and
    computational/latency considerations?
  \item How do temporal parameter choices (window length, stride,
    context length, preictal horizon) correlate with alarm-based
    performance and latency trade-offs in real-time use?
  \item What evidence exists for robustness to domain shift
    (patient/device/site/activity), including cross-dataset/external
    validation, and how do these factors affect alarm-based metrics?
  \item Which dataset gaps (class imbalance, limited prospective
    wearable data) are reported, and what is the measured impact of
    augmentation or synthetic data on alarm-based outcomes without
    introducing bias?
  \item How do personalization strategies (calibration, threshold
    adaptation, patient-specific fine-tuning) shift operating points
    (sensitivity vs FAR/h) under edge constraints?
  \item What alarm logic and thresholding designs (e.g., multi-window
    confirmation, refractory periods) are used to minimize FAR/h while
    preserving sensitivity and acceptable latency?
\end{enumerate}

\paragraph{Contribution}
The review will synthesize model design levers (representation, temporal
modeling, fusion, personalization, alarm logic) and evaluation
practices, clarifying how they influence alarm-based performance and
robustness, and deriving practical recommendations for achieving
clinically meaningful low FAR/h with high sensitivity in ambulatory
deployment.

\subsection{Methodology Overview and Structure}
The approach follows structured IS/medical literature guidance (concept
matrix inspired by \cite{websterAnalyzingPastDetermine2002}) and PRISMA
principles for transparent selection.

\paragraph{Scope}
Peer-reviewed human studies (journal articles, 2015--2025) using
non-EEG wearable/autonomic signals (ECG/HRV, PPG, EDA, ACC) for
seizure detection or preictal prediction.

\paragraph{Databases and Search Strategy}
Multi-stage search: (1) Scoping via Google Scholar to identify
terminology and candidate papers; (2) Formal queries in IEEE Xplore,
PubMed, and Scopus using controlled keyword strings ("epilepsy" AND
("seizure detection" OR "prediction") AND ("ECG" OR "HRV" OR
"PPG" OR "photoplethysmography" OR "EDA" OR "electrodermal" OR
"accelerometer" OR "wearable") AND ("deep learning" OR "machine
learning")) with deduplication; (3) Backward and forward citation
chasing on key seed papers.

\paragraph{Inclusion Criteria}
Human subject data; wearable or clinical acquisition of autonomic/
movement signals used singly or in multimodal fusion; seizure detection
or preictal prediction tasks; reported performance metrics (at least
sensitivity and a false-alarm measure) or sufficient methodological
detail for extraction.

\paragraph{Exclusion Criteria}
EEG-only or invasive studies; purely algorithmic or simulation-only work
without human data; animal-only studies; absence of seizure-related
task; no evaluable performance metrics; non-English/German language if
translation unavailable.

\paragraph{Evaluation Focus}
Primary: alarm-based sensitivity, FAR/h, latency. Secondary: AUC,
F1-score, specificity, PPV where available. Preictal horizons (e.g.,
5--60 min before onset) will be recorded when prediction tasks are
reported.

\paragraph{Screening Procedure}
Title/abstract screening followed by full-text eligibility assessment.
Reasons for exclusion documented. Counts summarized in a PRISMA-style
flow diagram (to be completed post-screening).

\paragraph{Data Extraction Dimensions}
A standardized template will capture: setting (clinical vs
ambulatory); modality (ECG/HRV, PPG, EDA, ACC); task (detection vs
prediction); model class (feature-based, 1D-CNN, LSTM, TCN,
transformer/attention); fusion strategy (early/feature,
intermediate/representation, late/decision); temporal parameters
(window length, stride, context length, preictal horizon); labeling
policy (ictal onset definition, preictal definition); alarm logic
(smoothing, multi-window confirmation, refractory period); personal-
ization (calibration, threshold adaptation, patient-specific
fine-tuning); study phase (retrospective, pseudo-prospective,
prospective); evaluation granularity (sample- vs alarm-based); metrics
(sensitivity, FAR/h, latency, AUC, F1, specificity, PPV); dataset
identity and availability; cohort size and class balance; interpreta-
bility/safety considerations; domain shift tests (cross-patient,
device/site/activity); external/cross-dataset validation; augmentation/
synthetic data usage; computational considerations (model size,
on-device feasibility, energy/latency notes).

\paragraph{Quality and Bias Flags}
Patient-preserving splits (e.g., LOSO); explicit external/cross-dataset
validation; alarm-based metric reporting; handling of class imbalance;
transparency of preictal horizon labels; prospective or pseudo-pros-
pective evaluation; model calibration/personalization procedure clarity.

\paragraph{Synthesis}
Narrative synthesis grouped by modality and task, integrating quantita-
tive ranges of alarm-based metrics and mapping them to modeling choices.
Comparative tables (concept matrix) will highlight patterns (e.g.,
transformer use vs improvement, fusion benefits under motion/noise,
personalization impact on FAR/h). Gaps (e.g., lack of prospective
multi-modality datasets) and future opportunities (compact attention
architectures, domain adaptation, uncertainty estimation) will be
identified.

\paragraph{Robustness and Generalization}
Special attention to domain shift: patient-level variability, device
placement differences (dominant wrist motion degrading wrist PPG/HR
accuracy), site/activity context, and cross-dataset transfers (e.g.,
TUH to RPAH/EPILEPSIAE) \parencite{yangMultimodalAISystem2022,andradePerformanceSeizurePrediction2024}.
Artifact and augmentation strategies (time-warping, jitter, scaling)
and out-of-distribution (OOD) detection adoption will be cataloged
\parencite{kalousiosECGbasedEpilepticSeizure2024,sethFeasibilityCardiacbasedSeizure2023}.

\paragraph{Healthcare-Specific Considerations}
Sensitivity vs FAR/h trade-off (alarm fatigue), interpretability
(clinician trust, safety mechanisms), latency constraints (early
intervention vs false positives), and regulatory suitability
(documented reasoning, fail-safes) are cross-cutting evaluation lenses.

\section{Technical Background (Clinical and Sensing Context)}
\subsection{Epileptic Seizures (Brief Overview)}
Seizures are classified by onset as focal (aware vs impaired awareness;
motor vs non-motor, possibly evolving to bilateral tonic--clonic) and
generalized (motor: tonic--clonic, myoclonic, atonic; non-motor:
absence) \parencite{fisherInstructionManualILAE2017}. Tonic--clonic
seizures commonly involve loss of consciousness and elevated
injury/SUDEP risk \parencite{hesdorfferCombinedAnalysisRisk2011}.

\subsection{Autonomic Markers and Wearable Sensing}
Sympathetic activation around seizures manifests as ictal tachycardia;
preictal heart-rate rise with reduced HRV (lower RMSSD/HF, elevated
LF/HF); increased EDA (tonic skin conductance level and phasic
responses); and limited ACC changes useful for artifact/context. PPG
yields pulse-rate variability analogous to HRV from ECG
\parencite{mironAutonomicBiosignalsSeizure2025,masonHeartRateVariability2024}.
Wearable heart-rate sensing aligns closely with chest-strap ECG under
motion (e.g., reported high concordance for consumer vs chest-strap
devices), though dominant-wrist motion degrades accuracy; chest-strap
ECG provides robust RR intervals during exercise
\parencite{wolfHearttoWearAssessingAccuracy2024a}.

\section{Modeling Landscape and Design Dimensions}
Classical pipelines (handcrafted HRV, EDA, ACC features with SVM, RF,
kNN) remain competitive in retrospective data
\parencite{sethFeasibilityCardiacbasedSeizure2023,masonHeartRateVariability2024}.
Deep temporal models (1D-CNN, LSTM, TCN) often improve phase
discrimination; attention and transformer variants are being explored
\parencite{yangMultimodalAISystem2022,pordoyEnhancedNonEEGMultimodal2025,abtahiIdentificationRelevantECG2025,kalousiosECGbasedEpilepticSeizure2024}.
Fusion across ECG, EDA, ACC, and PPG at feature, representation, or
decision level aims to exploit complementary autonomic and motion cues
under noise/non-stationarity \parencite{yangMultimodalAISystem2022,pordoyEnhancedNonEEGMultimodal2025}.
Personalization (patient-specific fine-tuning, threshold calibration)
can shift operating points toward lower FAR/h without substantial
sensitivity loss \parencite{andradePerformanceSeizurePrediction2024}.
Robustness evaluation via cross-dataset transfer emphasizes alarm-
over sample-based metrics \parencite{yangMultimodalAISystem2022,andradePerformanceSeizurePrediction2024}.
Artifact handling, augmentation (time-warping, jitter, scaling), and
OOD detection support generalization
\parencite{kalousiosECGbasedEpilepticSeizure2024,sethFeasibilityCardiacbasedSeizure2023}.
Deployment constraints (latency, energy, memory) motivate compact
sequence models and event-driven alarm logic with clear, explainable
rationales.

\section*{Provisional Table of Contents (\~15 Pages Text)}
\begin{enumerate}
  \item Introduction (Problem \& Motivation) -- 1.0 p
  \item Aim and Research Questions -- 0.5 p
  \item Methodology (search, selection, extraction, concept matrix,
    PRISMA overview) -- 1.0 p
  \item Technical Background (seizure taxonomy; autonomic markers;
    sensing accuracy) -- 1.0 p
  \item Modeling Landscape (features vs DL; CNN/LSTM/TCN;
    attention/transformers; fusion; personalization; robustness; alarm
    logic; deployment) -- 5--6 p
  \item Results (concept matrix; comparative ranges; modality/model
    comparison; temporal \& personalization effects) -- 3--4 p
  \item Discussion (trade-offs: sensitivity vs FAR/h; robustness;
    interpretability; deployment feasibility) -- 1.5--2 p
  \item Limitations and Future Work -- 0.5--1 p
  \item Conclusion -- 0.5 p
\end{enumerate}

\section*{Planned Output and Next Steps}
Next steps: finalize search strings; execute database queries; dedupli-
cate; complete PRISMA counts; populate extraction template; build
concept matrix; perform narrative synthesis and comparative analysis;
draft results and discussion sections.

\section*{Notation and Terminology Consistency}
Alarm-based metrics refer to event-level evaluation (alarms per hour,
latency, sensitivity). Sample-based metrics refer to per-window/predic-
tion classification accuracy. FAR/h = false alarms per hour. Patient-
preserving splits avoid leakage of subject-specific patterns.

\section*{Expected Contributions}
1) Consolidated mapping of modeling choices to alarm-based performance
for non-EEG wearable seizure tasks. 2) Identification of effective
personalization and fusion strategies under ambulatory constraints. 3)
Clarification of evaluation practices (retrospective vs pseudo-/prospec-
tive) and robustness gaps. 4) Recommendations for future architec-
tures (efficient attention, domain adaptation) and dataset needs
(prospective multimodal wearable corpora).

\section*{Limitations (Anticipated)}
Potential variability in metric definitions (e.g., differing preictal
horizons, alarm confirmation logic) may hinder strict quantitative
aggregation; reliance on published summaries may limit access to raw
performance distributions; scarcity of prospective transformer-based
wearable studies may constrain conclusions about attention benefits.

\section*{Ethical and Safety Considerations}
High sensitivity is critical for patient safety, but minimizing FAR/h
reduces burden and anxiety; interpretability and transparent alarm
rationale support clinician trust; personalization must avoid overfit-
ting that could degrade generalization in unseen contexts.

\section*{Conclusion (Preview)}
The review will integrate technical modeling trends in time-series deep
learning with clinical seizure monitoring needs, aiming to distill
design principles for low-burden, robust, and interpretable wearable
seizure detection and prediction systems.
